\documentclass[conference]{IEEEtran}
\usepackage{amsmath, amssymb, amsfonts}
\usepackage{graphicx}
\usepackage{algorithmic}
\usepackage{textcomp}
\usepackage{xcolor}

\begin{document}

\title{AI-Native Multicast Routing in FANETs using Graph Neural Networks and Deep Hierarchical Reinforcement Learning for Disaster Recovery Video Streaming}

\author{
\IEEEauthorblockN{Nitish M}
\IEEEauthorblockA{\textit{Computer Science Engineering} \\
\textit{St. Joseph's Institute of Technology}\\
Chennai, India \\
mnitish1705@gmail.com}
\and
\IEEEauthorblockN{Padmesh V A}
\IEEEauthorblockA{\textit{Computer Science Engineering} \\
\textit{St. Joseph's Institute of Technology}\\
Chennai, India \\
vapadmesh@gmail.com}
\and
\IEEEauthorblockN{Dr. V Sabaresan}
\IEEEauthorblockA{\textit{Computer Science Engineering} \\
\textit{St. Joseph's Institute of Technology}\\
Chennai, India}
}

\maketitle

\begin{abstract}
Flying Ad Hoc Networks (FANETs) present routing challenges due to mobility-driven topology changes and strict energy-delay constraints. This paper presents a Graph Neural Network (GNN) assisted Hierarchical Deep Reinforcement Learning (HDRL) framework for energy-aware multicast routing in FANETs. The implemented system uses two DQN-based controllers: a Meta controller for relay-node selection and an Intrinsic controller for hop-by-hop forwarding. The state encoder is a two-layer GCN that consumes 8-dimensional node features and topology adjacency, producing 16-dimensional embeddings. Path generation includes loop prevention through visited-node filtering and a strict maximum hop budget of 15. The reward function jointly optimizes packet delivery, latency, and energy usage with per-hop credit sharing during intrinsic learning. We implement and evaluate the approach in a Python simulation pipeline (Gym-style environment, 50 UAVs in a 2000x2000 $m^2$ region) with train/evaluate/demo/infer modes and AODV baseline comparison. The study establishes a reproducible algorithmic prototype and reports current performance characteristics and limitations for future improvement.
\end{abstract}

\begin{IEEEkeywords}
Flying Ad Hoc Networks, FANET, Multicast Routing, Hierarchical Reinforcement Learning, Deep Q-Network, Graph Neural Network, Energy-Aware Routing.
\end{IEEEkeywords}

\section{INTRODUCTION}
Flying Ad Hoc Networks (FANETs) formed by cooperating unmanned aerial vehicles (UAVs) are increasingly used in applications such as disaster response, environmental monitoring, precision agriculture, and tactical communications. Compared to ground MANETs, FANETs operate in three dimensions, have higher node speeds, and face more frequent topology changes, which significantly reduces route lifetime and complicates multicast distribution to multiple receivers. Standard MANET routing protocols (e.g., AODV, OLSR) and swarm-inspired optimizers (e.g., Artificial Bee Colony variants) can provide baseline functionality, but they are often brittle under extreme mobility and do not learn to trade off multiple competing QoS objectives online. Recent surveys highlight these FANET challenges and call for adaptive, learning-based routing solutions that exploit local observations to make robust forwarding decisions.

This work proposes an energy-aware Hierarchical Deep Reinforcement Learning (HDRL) framework empowered by Graph Neural Networks (GNN) for multicast routing in FANETs. Instead of static route discovery and periodic flooding, a two-level hierarchy learns routing behavior from graph-structured observations: a Meta controller selects relay candidates and an Intrinsic controller performs hop-by-hop forwarding. The implementation uses replay buffers of 10,000 and 20,000 transitions for the two controllers, respectively, and 16-dimensional GNN embeddings.

Our contributions are: (1) a hierarchical MDP-style formulation for multicast routing with a multi-objective reward balancing delivery, latency, and energy via hop-level reward sharing; (2) a two-layer GCN state encoder with 8 node features and link-aware topology adjacency; (3) an end-to-end HDRL training and inference pipeline for 50 UAVs with operational modes for training, evaluation, demo, and checkpointed inference; and (4) a reproducible simulation-based comparison against an AODV baseline with transparent reporting of current strengths and limitations.

\section{LITERATURE SURVEY}
Routing in FANETs has been studied extensively; surveys emphasize route instability due to high mobility, limited energy of onboard UAVs, and the need for three-dimensional mobility models and link quality estimation. Approaches range from topology-based protocols adapted from MANETs to hierarchical and geographic routing tailored to FANET use cases. Swarm intelligence and bio-inspired metaheuristics (e.g., Artificial Bee Colony and its adaptive variants) have been applied to MANET routing to optimize path selection with objectives such as hop count and residual energy. These methods produce good offline or slowly changing solutions but require frequent re-optimization in rapidly changing networks, which undermines responsiveness in FANETs.

Reinforcement learning (RL), including multi-agent RL variants, has recently been explored for networking tasks where dynamic adaptation is advantageous. Prior works demonstrate RL for offloading decisions, traffic scheduling, and delay-constrained resource allocation. In this paper, we focus on a hierarchical DQN routing design rather than a full multi-agent actor-critic setup. This design choice keeps the action space tractable while retaining adaptivity for dynamic FANET topologies and energy/latency tradeoffs.

\section{SYSTEM ARCHITECTURE}
We design a layered, modular architecture that separates simulation/control, agent learning, and runtime packet-forwarding components to enable both research and prototyping (Fig. 1 conceptual).

\subsection{Network and Simulation Layer}
Defines the FANET network behavior used during training and evaluation. The environment simulates 50 UAVs spread across a 2000x2000 $m^2$ area with a multicast group of 1 source to 5 destinations. Mobility follows Random Waypoint behavior with dynamic altitude state and bounded speed. The link model computes signal strength, delay, loss, and bandwidth surrogates from inter-node distance. This layer is implemented as a Python Gym-style environment; an optional NS-3 bridge exists but is not required in the current experiments.

\subsection{Meta Controller Layer (Global Strategy)}
A high-level Meta Controller acts as the global strategist. By observing aggregated network features through GNN embeddings, it determines strategic relay nodes or high-level milestones for multicast propagation, buffering 10,000 transitions to stabilize strategic policies.

\subsection{Intrinsic Controller Layer (Local Execution)}
A lower-level Intrinsic Controller manages the local, hop-by-hop execution of paths bounded by the Meta Controller's targets. It ensures valid packet forwarding, avoids routing loops via explicit visited sets, strictly enforces a 15-hop maximum path limit, and stabilizes learning across a 20,000-transition replay buffer.

\subsection{Routing \& Maintenance Layer}
Contains explicit logic to translate hierarchical controller actions into viable network topologies, including loop avoidance, graph edge validations, and integration with the operational modes (\texttt{train}, \texttt{evaluate}, \texttt{demo}, \texttt{infer}) for checkpoint evaluation.

\subsection{Monitoring \& Replay Layer}
Collects trace logs (packet deliveries, latencies, and energy consumption) for offline analysis and model comparison across checkpoints and baselines.

\section{PROPOSED METHODOLOGY}
We model routing as a hierarchical decision process with two DQN policies operating on graph-encoded state. Below we detail state, action, reward, training algorithm, and route maintenance.

\subsection{Problem Formulation (Hierarchical MDP)}
At each simulation step, the environment provides a graph snapshot, node-feature matrix, adjacency matrix, source node, and destination set. A GCN encoder maps this input to node embeddings. The Meta controller then selects relay candidates from available nodes, and the Intrinsic controller constructs per-destination paths by selecting next hops from current neighbors.

\textbf{State}
The Meta state is formed by concatenating source embedding, mean destination embedding, and global mean embedding. The Intrinsic state concatenates current-node embedding, destination embedding, and mean neighbor embedding.

\textbf{Actions}
\begin{itemize}
    \item Meta action: choose relay-node indices from a candidate list.
    \item Intrinsic action: choose the next-hop neighbor for each step in a source-to-destination path.
\end{itemize}

\textbf{Transition and Constraints}
Paths are built with visited-node filtering and a strict hop cap of 15 to prevent loops and runaway forwarding. Episodes terminate on step budget exhaustion or insufficient live nodes.

\subsection{State Encoding (GNN Embedding)}
Since FANETs have a dynamic topology with varying numbers of neighboring nodes, local observations must be encoded efficiently into a fixed dimensionality. The environment models an 8-dimensional node feature state (including position, speed, energy, current traffic load, and queue size) and a 4-dimensional edge relationship between neighboring UAVs (signal strength, delay, loss, bandwidth).

For each global graph snapshot, a Graph Convolutional Network (GCN) is employed with two layers. The node features $x_i$ and edge features $e_{ij}$ flow as:
$$ x_i^{l+1} = \sigma \left( \sum_{j \in \mathcal{N}(i)} \frac{1}{c_{ij}} W^{(l)} x_j^{(l)} \right) $$

The 2-layer GCN projects the variable-length topology down to a 32-dimensional intermediate layer, and finally a 16-dimensional embedding layer. This learned 16D embedding compactly represents the local topology of 50 UAV nodes.

\subsection{Reward Design}
To encourage efficiency while avoiding routing loops, the environment explicitly computes a holistic objective bridging delivery, path optimization, latency, and longevity. The reward formula balances multiple goals as follows:
$$ r(t) = 10 \times D - 5 \times (P - D) - 2 \times \text{delay} - 1 \times E_{used} + 2 \times S $$

Where:
\begin{itemize}
    \item $D$ is the count of packets successfully delivered to their destinations.
    \item $P - D$ represents undelivered paths acting as redundant traversals.
    \item $\text{delay}$ measures the hop-to-hop latency accumulated over the route.
    \item $E_{used}$ is the total residual energy spent.
    \item $S = \text{alive\_nodes} / \text{total\_nodes}$ explicitly encourages the survival of nodes over time.
\end{itemize}
Furthermore, rather than evaluating endpoints, we adopt explicit \textit{hop-level reward sharing}, giving each intermediary the proportionate responsibility $r\_i(t) = r(t) / h$, where $h$ is the number of hops. A strict 15-hop limit caps the maximum routing traversals before packets drop, alongside \textit{visited sets} explicitly avoiding cyclical retransmissions.

\subsection{Algorithm and Training (Hierarchical DQN)}
We employ a two-tier Hierarchical Reinforcement Learning framework with Deep Q-Networks (DQN). The Meta controller manages strategic relay assignments with a replay buffer of 10,000 transitions. The Intrinsic controller handles path construction with a replay buffer of 20,000 transitions. Simulation is automated with periodic checkpointing and checkpoint recovery for \texttt{infer} mode.

\subsection{Multicast Path Construction \& Loop Prevention}
The current implementation realizes application-layer multicast as a set of per-destination forwarding paths. Visited-node constraints and the 15-hop limit prevent cycles and excessive route expansion.

\subsection{Scalability \& Transfer}
To scale to large swarms (50+ UAVs) distributed across vast terrains, we implicitly reuse the permutation-invariant representations encoded natively by the GNNs. 

\section{SYSTEM IMPLEMENTATION}
\subsection{Development Environment}
\begin{itemize}
    \item \textbf{Simulation:} Python Gym-style FANET simulator with 50 UAVs over a 2000x2000 $m^2$ region.
    \item \textbf{ML stack:} PyTorch Hierarchical DQN controllers with a two-layer GCN encoder.
    \item \textbf{Mobility:} Random Waypoint motion with altitude state and maximum speed 15 m/s.
    \item \textbf{Operational modes:} \texttt{--mode train / evaluate / demo / infer}, including checkpoint-based inference and plot generation.
\end{itemize}

\begin{table}[htbp]
\centering
\caption{Simulation Environment and Hyperparameter Settings}
\renewcommand{\arraystretch}{1.3}
\begin{tabular}{|l|l|}
\hline
\textbf{Parameter} & \textbf{Value} \\
\hline
\textbf{Network and Simulation Settings} & \\
Simulation Area & $2000 \times 2000$ $m^2$ \\
Total UAV Nodes & 50 \\
Multicast Group Setup & 1 Source, 5 Destinations \\
Mobility Model & Random Waypoint (with altitude state) \\
Maximum Node Speed & 15 m/s \\
Link Model & Distance-based signal/delay/loss/bandwidth surrogate \\
\hline
\textbf{Deep RL and Agent Settings} & \\
GNN Node Features & 8 (pos, speed, energy, load, queue) \\
GNN Edge Features & 4 (signal, delay, loss, bw) \\
GNN Embeddings & 16-Dimensional \\
Hidden Units (MLPs) & 128 \\
Meta Controller Buffer & 10,000 transitions \\
Intrinsic Controller Buffer & 20,000 transitions \\
Learning Rate $\alpha$ & $1 \times 10^{-3}$ (Adam) \\
Discount Factor $\gamma$ & 0.99 \\
Maximum Hops Allowed & 15 \\
\hline
\end{tabular}
\label{tab:sim_parameters}
\end{table}

\subsection{Agent Architecture}
\textbf{Actors:} Hierarchical Controllers consuming an integrated 16-Dimensional GNN summary.
\textbf{Q-Network:} State vectors use $3 \times 16 = 48$ dimensions, with MLP hidden layers of width 128.

\subsection{Training Details}
\begin{itemize}
    \item \textbf{Optimization:} Adam with Meta $(lr = 1e-3)$ and Intrinsic $(lr = 1e-3)$ decoupled models.
    \item \textbf{Discount factor:} $\gamma = 0.99$.
    \item \textbf{Memory Storage:} Replay buffers of 10,000 (Meta) and 20,000 (Intrinsic) transitions.
    \item \textbf{Checkpoints:} Systematically captured and auto-recovered by history heuristics logging $training\_history.json$.
\end{itemize}

\subsection{Baselines and Comparison Setup}
We compare GNN-HDRL against a shortest-path AODV-style baseline (and weighted shortest-path variant in code) using packet delivery ratio (PDR), average delay, and energy-related metrics.

\begin{figure}[htbp]
    \centering
    % Placeholder for the provided flowchart image
    \includegraphics[width=\linewidth]{system_flowchart.png}
    \caption{System Architecture Flowchart}
    \label{fig:flowchart}
\end{figure}

\section{EXPERIMENTAL RESULTS}
The proposed GNN-HDRL framework was evaluated in the implemented Python simulation setup with 50 UAVs and multicast sessions (1 source, 5 destinations). We report reproducible artifact-driven observations from saved training and inference histories.

\subsection{Quantitative Performance Evaluation}
For the current saved run (100 inference episodes), summary statistics from \texttt{inference\_history.json} are: mean PDR $\approx 0.526$, mean delay $\approx 0.0239$ s, and mean episode energy usage $\approx 5.35$ (simulator units). Across 1000 recorded training episodes, mean PDR is $\approx 0.521$ with corresponding delay/energy traces stored in \texttt{training\_history.json}. These results indicate stable training behavior but also room for policy improvement.

\begin{figure}[htbp]
    \centering
    % Placeholder for PDR results
    \includegraphics[width=\linewidth]{pdr_comparison.png}
    \caption{Packet Delivery Ratio (PDR) Performance Comparison}
    \label{fig:pdr}
\end{figure}

\textbf{GNN-HDRL Routing (Current Run)}
\begin{itemize}
    \item Mean PDR around 0.53 over 100 inference episodes.
    \item Mean end-to-end delay around 23.9 ms in simulator time.
\end{itemize}

\begin{figure}[htbp]
    \centering
    % Placeholder for delay results
    \includegraphics[width=\linewidth]{latency_comparison.png}
    \caption{Average End-to-End Delay Comparison}
    \label{fig:delay}
\end{figure}

\textbf{AODV-Style Baseline Routing}
\begin{itemize}
    \item Used as a shortest-path reference for comparison.
    \item In this project stage, baseline and HDRL performance are scenario-dependent; further tuning is required for consistent HDRL gains.
\end{itemize}

\begin{table}[htbp]
\centering
\caption{Comparison Matrix: GNN-HDRL (Current Implementation) vs. AODV-Style Baseline}
\renewcommand{\arraystretch}{1.3}
\begin{tabular}{|l|c|c|}
\hline
\textbf{Metric / Feature} & \textbf{GNN-HDRL (Our Model)} & \textbf{AODV (Baseline)} \\
\hline
\textbf{Routing Strategy} & GNN + Hierarchical DQN & Shortest-Path Routing \\
\textbf{Packet Delivery Ratio} & Mean $\approx 0.53$ (inference run) & Scenario-dependent \\
\textbf{Average Latency} & Mean $\approx 23.9$ ms (simulator) & Scenario-dependent \\
\textbf{Energy Awareness} & Included in reward function & Not learning-based \\
\textbf{Adaptability} & Learns from replay experience & Rule-based \\
\hline
\end{tabular}
\label{tab:comparison_matrix}
\end{table}

\begin{figure}[htbp]
    \centering
    % Placeholder for energy/lifetime results
    \includegraphics[width=\linewidth]{energy_consumption.png}
    \caption{Energy Consumption and Network Lifetime}
    \label{fig:energy}
\end{figure}

The experimental results show that the implemented GNN-HDRL framework is functional and reproducible, with measurable delivery, delay, and energy behavior under dynamic topology. At the current tuning stage, results should be interpreted as a proof-of-concept rather than a definitive replacement for classical baselines in all scenarios.

\subsection{Behavioral Analysis of the Learned Policy}
Policy behavior indicates that the intrinsic controller prefers feasible neighbor expansions constrained by connectivity and visited-node filtering. Because action selection is learned from replayed transitions and reward shaping, route choices are influenced by delivery success, latency penalties, and energy penalties rather than hop count alone.

\subsection{Ablation Study (Planned)}
Formal ablation experiments (e.g., removing energy penalty or replacing GNN embeddings with fixed local features) are planned for future work. The current repository does not yet include dedicated ablation scripts or isolated checkpoints for these variants.

\subsection{Limitations of the Simulation Study}
Although the simulation results provide useful insight, several limitations remain. First, the experiments were conducted in the built-in Python simulator; optional NS-3 bridging was not used for the reported results. Real UAV networks may exhibit additional radio and hardware effects not captured here.

Second, outcomes depend on hyperparameters (exploration schedule, reward coefficients, replay dynamics) and mobility assumptions. Finally, the present implementation performs application-layer multicast through per-destination path construction, not a strict shared-tree network-layer multicast protocol. These limitations define clear directions for improvement.

\section{CONCLUSION}
We presented a GNN-assisted Hierarchical DQN framework for energy-aware multicast routing in FANET simulation. The current implementation combines graph-based state encoding, hierarchical relay/next-hop policies, loop-safe path construction, and checkpointed train/infer workflows. Results from saved runs demonstrate a functional and reproducible prototype, while also highlighting that additional tuning and broader benchmarking are needed for consistent superiority over classical baselines.

\section{FUTURE DEVELOPMENT}
\begin{itemize}
    \item \textbf{Real-world prototyping:} Port the learned policies to onboard UAV hardware (edge inference) and validate over flight tests with real radios.
    \item \textbf{Hybrid learning:} Combine supervised bootstrapping from classical routing heuristics with RL fine-tuning to reduce initial exploration and speed safe deployment.
    \item \textbf{Mobility-aware collaboration:} Extend actions to include small collaborative mobility adjustments (positioning assistance) to proactively improve connectivity for multicast.
    \item \textbf{End-to-end representation learning:} Train and checkpoint the GNN encoder jointly with the controllers instead of using it as a fixed feature extractor.
    \item \textbf{True multicast trees:} Add shared-tree construction/enforcement so relay decisions are strictly reflected in final forwarding structures.
    \item \textbf{Security and robustness:} Investigate adversarial and Byzantine-tolerant training to mitigate routing attacks and misbehaving nodes.
    \item \textbf{Energy harvesting models:} Incorporate energy harvesting (solar) models for long-endurance UAVs and adapt policies accordingly.
\end{itemize}

\begin{thebibliography}{00}
\bibitem{b1} F. Pasandideh, M. Naderan, and M. R. Javan, ``A Review of Flying Ad Hoc Networks: Key Characteristics, Applications and Wireless Technologies,'' Remote Sensing, vol. 14, no. 18, pp. 1--30, 2022.
\bibitem{b2} M. F. Khan, S. A. Hassan, H. Jung, and M. S. Hossain, ``Routing Schemes in Flying Ad Hoc Networks: A Survey,'' Sensors, vol. 19, no. 1, pp. 1--29, 2019.
\bibitem{b3} Y. Zhang, J. Liu, and H. Li, ``MANIA-F: Multi-Agent Intra-FANET Framework Based on Multi-Agent Proximal Policy Optimization,'' IEEE Access, vol. 10, pp. 82345--82357, 2022.
\bibitem{b4} H. Wang, X. Zhang, and Y. Chen, ``Multi-Agent Deep Reinforcement Learning-Based Fine-Grained Traffic Scheduling,'' Future Internet, vol. 16, no. 4, pp. 1--17, 2024.
\bibitem{b5} D. Karaboga and B. Basturk, ``Artificial Bee Colony (ABC) Optimization Algorithm for Solving Constrained Optimization Problems,'' Applied Soft Computing, vol. 8, no. 1, pp. 687--697, 2008.
\bibitem{b6} S. Hayat, E. Yanmaz, and C. Bettstetter, ``Experimental Analysis of Multipoint-to-Point UAV Communications with IEEE 802.11n and 802.11ac,'' IEEE Communications Letters, vol. 19, no. 12, pp. 2246--2249, 2015.
\bibitem{b7} Y. Zeng, R. Zhang, and T. J. Lim, ``Wireless Communications with Unmanned Aerial Vehicles: Opportunities and Challenges,'' IEEE Communications Magazine, vol. 54, no. 5, pp. 36--42, 2016.
\bibitem{b8} L. Gupta, R. Jain, and G. Vaszkun, ``Survey of Important Issues in UAV Communication Networks,'' IEEE Communications Surveys \& Tutorials, vol. 18, no. 2, pp. 1123--1152, 2016.
\bibitem{b9} J. Chen, S. Yin, and X. Zhang, ``Reinforcement Learning for Routing in Dynamic Networks,'' IEEE Transactions on Network Science and Engineering, vol. 7, no. 3, pp. 1416--1427, 2020.
\bibitem{b10} R. Lowe, Y. Wu, A. Tamar, J. Harb, P. Abbeel, and I. Mordatch, ``Multi-Agent Actor-Critic for Mixed Cooperative-Competitive Environments,'' Advances in Neural Information Processing Systems (NeurIPS), pp. 6379--6390, 2017.
\bibitem{b11} J. Schulman, F. Wolski, P. Dhariwal, A. Radford, and O. Klimov, ``Proximal Policy Optimization Algorithms,'' arXiv preprint arXiv:1707.06347, 2017.
\bibitem{b12} P. Sunehag, G. Lever, A. Gruslys, et al., ``Value-Decomposition Networks for Cooperative Multi-Agent Learning,'' Proceedings of the International Conference on Autonomous Agents and Multiagent Systems (AAMAS), pp. 2085--2087, 2018.
\bibitem{b13} T. Clausen and P. Jacquet, ``Optimized Link State Routing Protocol (OLSR),'' IETF RFC 3626, 2003.
\bibitem{b14} C. Perkins and E. Royer, ``Ad-hoc On-Demand Distance Vector Routing,'' Proceedings of IEEE WMCSA, pp. 90--100, 1999.
\bibitem{b15} S. Park and D. Kim, ``Energy Efficient Routing Protocols for UAV Networks,'' IEEE Access, vol. 8, pp. 19734--19744, 2020.
\bibitem{b16} X. Liu, Z. Li, and J. Zhao, ``Graph Neural Network Based Routing in Dynamic Wireless Networks,'' IEEE Transactions on Mobile Computing, vol. 21, no. 5, pp. 1823--1836, 2022.
\bibitem{b17} K. Zhang, Z. Yang, and T. Basar, ``Multi-Agent Reinforcement Learning: A Selective Overview of Theories and Algorithms,'' Handbook of Reinforcement Learning and Control, pp. 321--384, 2021.
\end{thebibliography}

\end{document}